% c06 —— §9 末（p022 上）+ §10 可逆多项式（p022 下–p026），pages p022–p026
% tag 共 41 条：(9.39)–(9.41)、(10.1)–(10.38)
% 注：tags_bypage.txt 中 p024 的 (10.3) 是行文引用造成的假 tag，本块只在 p023 写出 (10.3) 一次。

而 $l\neq 0$ 的项则为
\begin{align}
    \sum_{l=1}^{3}\left[(t\bar t)^{\frac{1}{2}-\frac{1}{4}}\left(\frac{t}{\bar t}\right)^{\frac{1}{4}l-\left[\frac{1}{4}l\right]-\frac{1}{2}}\right]^{4}
    &= \sum_{l=1}^{3}(t\bar t)^{2-1}\left(\frac{t}{\bar t}\right)^{l-2} \tag{9.39}\\
    &= \sum_{l=1}^{3}t^{l-1}\bar t^{3-l} \tag{9.40}\\
    &= \sum_{p=0}^{2}t^{p}\bar t^{2-p} \tag{9.41}
\end{align}

用 Vafa 公式 (9.26)，对与 7555 个权相对应的三维 Calabi–Yau 流形画出期待中的 Hodge 数，就会发现其中能找到许多拓扑镜像对；这一事实当时令代数几何学家惊讶不已。请看 [CLS90, Fig. 1] 中那幅令人印象深刻的图。另一方面，7555 个之中的 850 个在 7555 个之中并不持有拓扑镜像。因此，如何在数学上理解 Vafa 公式，以及为剩下的 850 个找出拓扑镜像，就成为接下来的问题。作为处理这些问题的方法，最有名的是使用环面几何的方法，但在讨论它之前，先介绍准备更少的、使用可逆多项式的方法。
% [Tutor Note: 原文作「7555 個の重み付きに対応する」，「付き」疑为衍字；译文按「7555 個の重みに対応する」（与 7555 个权相对应）处理。]

\section{可逆多项式}

称满足下列条件的多项式 $f\in\C[x_1,\ldots,x_n]$ 为可逆多项式（invertible polynomial）：

\begin{enumerate}
\item 存在以非负整数为元素的可逆矩阵 $A=(a_{ij})_{i,j=1}^{n}$，使得
\begin{equation}
    f=\sum_{i=1}^{n}\prod_{j=1}^{n}x_j^{a_{ij}} \tag{10.1}
\end{equation}
成立。
\item $f$ 在原点有孤立临界点。
\end{enumerate}

在原点有孤立临界点常常被称为非退化（non-degenerate）（例如见 [AGZV85, p. 192]）。注意这与 Morse 理论中临界点的非退化性不同。

对多项式 $f(x_1,\ldots,x_m)\in\C[x_1,\ldots,x_m]$ 与 $g(y_1,\ldots,y_n)\in\C[y_1,\ldots,y_n]$，$f$ 与 $g$ 的 Sebastiani–Thom 和（Sebastiani–Thom sum）由
\begin{equation}
    f(x_1,\ldots,x_m)+g(y_1,\ldots,y_n)\in\C[x_1,\ldots,x_m,y_1,\ldots,y_n] \tag{10.2}
\end{equation}
定义。

\begin{exercise}
证明：任意可逆多项式都是下列三种型式的多项式的 Sebastiani–Thom 和（提示：见 [KS92b]）：
\begin{itemize}
\item Fermat 型（Fermat type）：$x^{p}$。
\item 链型（chain type）：$x_1^{p_1}x_2+x_2^{p_2}x_3+\cdots+x_{n-1}^{p_{n-1}}x_n+x_n^{p_n}$。
\item 环型（loop type）：$x_1^{p_1}x_2+x_2^{p_2}x_3+\cdots+x_{n-1}^{p_{n-1}}x_n+x_n^{p_n}x_1$。
\end{itemize}
\end{exercise}

若 $f$ 是可逆多项式，则由
\begin{equation}
    \check f=\sum_{i=1}^{n}\prod_{j=1}^{n}x_j^{a_{ji}} \tag{10.3}
\end{equation}
定义的 $f$ 的转置多项式（Berglund–Hübsch transpose）也是可逆的。

由 $A$ 是可逆矩阵这一条件，任意可逆多项式都是加权齐次多项式，且对应的权除常数倍外唯一确定。按最大公约数为 1 这一条件固定常数倍的自由度，所得的权记为
\begin{equation}
    (a_1,\ldots,a_n;h):=\deg(x_1,\ldots,x_n;f) \tag{10.4}
\end{equation}
并称为既约权系（reduced weight system）。

对给定的既约权系，对应的可逆多项式即使存在也不唯一确定。例如，可逆多项式
\begin{equation}
    \begin{array}{r|l}
        4F_1 & x^4+y^4+z^4+w^4 \\
        2F_1+C_2 & x^4+y^4+z^3w+w^4 \\
        2F_1+L_2 & x^4+y^4+z^3w+w^3z \\
        F_1+C_3 & x^4+y^3z+z^3w+w^4 \\
        F_1+L_3 & x^4+y^3z+z^3w+w^3y \\
        2C_2 & x^3y+y^4+z^3w+w^4 \\
        C_2+L_2 & x^3y+y^4+z^3w+w^3z \\
        2L_2 & x^3y+y^3z+z^3w+w^3z \\
        C_4 & x^3y+y^3z+z^3w+w^4 \\
        L_4 & x^3y+y^3z+z^3w+w^3x
    \end{array}
    \tag{10.5}
\end{equation}
对应的既约权系全都是 $(1,1,1,1;4)$。

另外，对定义拟光滑的加权 Calabi–Yau 超曲面的权，对应的可逆多项式也未必一定存在。例如，给出 K3 曲面的 95 个权之一 $(2,4,5,11;22)$，可用于构成可逆多项式的单项式只有
\begin{equation}
    x^{11},\ x^{9}y,\ xy^{5},\ xz^{4},\ w^{2} \tag{10.6}
\end{equation}
这 5 个，而用它们构不成可逆多项式，因此可知与该权对应的可逆多项式不存在。对 95 个权中的每一个作同样的考察，便知 95 个之中有 91 个持有对应的可逆多项式，但
\begin{equation}
    \begin{array}{c|l}
        \mathrm{No.} & \text{权} \\ \hline
        75 & (2,4,5,11;22) \\
        90 & (4,6,7,17;34) \\
        91 & (5,6,8,19;38) \\
        93 & (3,4,10,17;34)
    \end{array}
    \tag{10.7}
\end{equation}
这 4 个不持有对应的可逆多项式\footnote{据与桥本健治、三浦真人的讨论。}。

可逆多项式 (10.1) 是加权齐次多项式，因此它定义对应的加权射影空间中的超曲面
\begin{equation}
    Y:=\bigl\{[x_1:\cdots:x_n]\in\PP(a_1,\ldots,a_n)\ \big|\ f(x_1,\ldots,x_n)=0\bigr\} \tag{10.8}
\end{equation}
由可逆多项式只在原点有临界点可知，$Y$ 是拟光滑的。$Y$ 成为（一般带有 Gorenstein 商奇点的）Calabi–Yau 流形的充分必要条件是
\begin{equation}
    a_1+\cdots+a_n=h \tag{10.9}
\end{equation}

对可逆多项式 (10.1)，其极大对角对称群（the group of maximal diagonal symmetries）由
\begin{equation}
    G_{\max}:=\left\{(\alpha_1,\ldots,\alpha_n)\in(\mathbb{G}_m)^{n}\ \middle|\ \prod_{j=1}^{n}\alpha_j^{a_{ij}}=1\ \ (i=1,\ldots,n)\right\} \tag{10.10}
\end{equation}
定义。
% [Tutor Note: 原文把条件「任意の i=1,…,n に対して」写在集合内部；为守「公式中不出现中文」的规范，译文以 (i=1,\ldots,n) 表示，含义不变。]

通过到对角子群的嵌入
\begin{equation}
    (\mathbb{G}_m)^{n}\to\mathrm{GL}_n(\C),\qquad (\alpha_1,\ldots,\alpha_n)\mapsto\operatorname{diag}(\alpha_1,\ldots,\alpha_n) \tag{10.11}
\end{equation}
极大对角对称群与 $\mathrm{GL}_n(\C)$ 的子群等同。同样地，对转置多项式 (10.3)，其极大对角对称群由
\begin{equation}
    \check G_{\max}:=\left\{(\alpha_1,\ldots,\alpha_n)\in(\mathbb{G}_m)^{n}\ \middle|\ \prod_{j=1}^{n}\alpha_j^{a_{ji}}=1\ \ (i=1,\ldots,n)\right\} \tag{10.12}
\end{equation}
定义。设 $\check f$ 所定义的加权射影空间中的超曲面为 $\check Y$，则 $\check G_{\max}$ 自然地作用于 $\check Y$。为使商空间再次成为（一般带有 Gorenstein 商奇点的）Calabi–Yau 流形，考虑的不是 $\check G_{\max}$ 本身，而是其子群
\begin{equation}
    \check G:=\check G_{\max}\cap\mathrm{SL}_n(\C) \tag{10.13}
\end{equation}
朴素地想说 $\check Y/\check G$ 是 $Y$ 的镜像，但 $Y$ 与 $\check Y/\check G$ 一般带有 Gorenstein 商奇点，故代之以考虑对应的光滑轨形\footnote{流形以欧氏空间为局部模型，而把「欧氏空间与作用于其上的群」这一组作为局部模型的流形之推广概念，就是轨形。例如，用 $\Z/2\Z$ 的作用 $z\mapsto -z$ 去除 $\C$ 所得的空间，作为拓扑空间与 $\C$ 同胚，但作为轨形并不同构。又，在代数几何上，轨形被作为光滑的 Deligne–Mumford 栈（smooth Deligne–Mumford stack）处理。关于栈的入门文献有例如 [Fan01, Góm01, Vis05]，教科书有 [LMB00, Ols16] 等。此外还有称为 Stacks project 的开源文献，作为参考文献非常方便，但篇幅庞大，通读实际上不可能。}\footnote{商奇点具有 crepant 消解时，也可以考虑 crepant 消解来代替对应的轨形。3 维及以下 Gorenstein 商奇点一定具有 crepant 消解，但 4 维及以上则未必存在。另外，对 2 维 Gorenstein 奇点，crepant 消解与极小奇点消解（即例外除子中不含自交数为 $-1$ 的曲线的奇点消解）一致，因而唯一；但 3 维及以上 crepant 消解即使存在也不一定唯一。这可以解释为一个 Calabi–Yau 流形存在多个镜像，也可以解释为双有理等价的 Calabi–Yau 流形不应加以区分。}。下面把这一点准确表述出来。

任意可逆多项式不仅是加权齐次多项式，而且是关于由 $n+1$ 个元素 $\vec x_1,\ldots,\vec x_n,\vec c$ 生成、具有 $n$ 个关系式
\begin{equation}
    a_{i1}\vec x_1+\cdots+a_{in}\vec x_n=\vec c,\qquad i=1,\ldots,n \tag{10.14}
\end{equation}
的秩 1 Abel 群 $L$ 所给的多项式环 $\C[x_1,\ldots,x_n]$ 的分次
\begin{equation}
    \deg x_i=\vec x_i,\qquad i=1,\ldots,n \tag{10.15}
\end{equation}
的次数为 $\vec c$ 的齐次多项式。这个 $L$ 是
\begin{equation}
    K:=\left\{(\alpha_1,\ldots,\alpha_n)\in(\mathbb{G}_m)^{n}\ \middle|\ \alpha_1^{a_{11}}\cdots\alpha_n^{a_{1n}}=\cdots=\alpha_1^{a_{n1}}\cdots\alpha_n^{a_{nn}}\right\} \tag{10.16}
\end{equation}
所定义的群的特征标群，而映射
\begin{equation}
    \begin{array}{ccc}
        K & \longrightarrow & \mathbb{G}_m \\[2pt]
        \text{\rotatebox[origin=c]{-90}{$\mapsto$}} & & \text{\rotatebox[origin=c]{-90}{$\mapsto$}} \\[2pt]
        (\alpha_1,\ldots,\alpha_n) & \longmapsto & \alpha_1\cdots\alpha_n
    \end{array}
    \tag{10.17}
\end{equation}
的核与极大对角对称群 $G_{\max}$ 自然地等同；
\begin{equation}
    1\to G_{\max}\to K\to\mathbb{G}_m\to 1 \tag{10.18}
\end{equation}
% [Tutor Note: 原文「多項式環の C[x1,…,xn]」中的「の」疑为衍字，译文按「多项式环 C[x1,…,xn]」处理。]

把 $G_{\max}$ 的特征标群记为
\begin{equation}
    G^{\vee}_{\max}:=\Hom(G_{\max},\mathbb{G}_m) \tag{10.19}
\end{equation}
由正合列 (10.18) 便得到
\begin{equation}
    1\to\Z\to L\to G^{\vee}_{\max}\to 1 \tag{10.20}
\end{equation}
由于 $G_{\max}$ 是有限 Abel 群，$G_{\max}$ 与 $G^{\vee}_{\max}$（非典范地）同构。把 $A$ 的逆矩阵写成
\begin{equation}
    A^{-1}=
    \begin{pmatrix}
        \varphi_1^{(1)} & \varphi_1^{(2)} & \cdots & \varphi_1^{(n)} \\
        \varphi_2^{(1)} & \varphi_2^{(2)} & \cdots & \varphi_2^{(n)} \\
        \vdots & \vdots & \ddots & \vdots \\
        \varphi_n^{(1)} & \varphi_n^{(2)} & \cdots & \varphi_n^{(n)}
    \end{pmatrix} \tag{10.21}
\end{equation}
则 $G_{\max}$ 由
\begin{equation}
    \rho_k=\left(\exp\left(2\pi\sqrt{-1}\varphi_1^{(k)}\right),\ldots,\exp\left(2\pi\sqrt{-1}\varphi_n^{(k)}\right)\right),\qquad k=1,\ldots,n \tag{10.22}
\end{equation}
生成。令
\begin{equation}
    \varphi_i=\varphi_i^{(1)}+\cdots+\varphi_i^{(n)},\qquad i=1,\ldots,n, \tag{10.23}
\end{equation}
并按
\begin{equation}
    h:=\min\left\{k\in\Z^{>0}\ \middle|\ k\varphi_i\in\Z\ \ (i=1,\ldots,n)\right\}, \tag{10.24}
\end{equation}
\begin{equation}
    a_i:=h\varphi_i,\qquad i=1,\ldots,n \tag{10.25}
\end{equation}
定义 $(a_1,\ldots,a_n;h)$，则它是使多项式 $f$ 加权齐次的既约权系。由于 $f$ 在原点有孤立临界点，
\begin{equation}
    \tilde Y:=\Spec\C[x_1,\ldots,x_n]/(f) \tag{10.26}
\end{equation}
是仅在原点有奇点的代数簇，而商栈
\begin{equation}
    \mathbb{Y}:=\left[\left(\tilde Y\setminus\mathbf{0}\right)/\mathbb{G}_m\right] \tag{10.27}
\end{equation}
是光滑轨形。

同样地，$f$ 的转置 $\check f$ 的极大对角对称性群 $\check G_{\max}$ 的生成元，可用 $\check A^{-1}$ 的列向量 $(\check\varphi_i^{(k)})_{i=1}^{n}$（它也是 $A^{-1}$ 的行向量）表示为
\begin{equation}
    \check\rho_k=\left(\exp\left(2\pi\sqrt{-1}\check\varphi_1^{(k)}\right),\ldots,\exp\left(2\pi\sqrt{-1}\check\varphi_n^{(k)}\right)\right),\qquad k=1,\ldots,n \tag{10.28}
\end{equation}
令
\begin{equation}
    \check\varphi_i=\check\varphi_i^{(1)}+\cdots+\check\varphi_i^{(n)},\qquad i=1,\ldots,n, \tag{10.29}
\end{equation}
并按
\begin{equation}
    \check h:=\min\left\{k\in\Z^{>0}\ \middle|\ k\check\varphi_i\in\Z\ \ (i=1,\ldots,n)\right\}, \tag{10.30}
\end{equation}
\begin{equation}
    \check a_i:=\check h\check\varphi_i,\qquad i=1,\ldots,n \tag{10.31}
\end{equation}
定义 $(\check a_1,\ldots,\check a_n;\check h)$，则它是使多项式 $\check f$ 加权齐次的既约权系。定义同态 $\check\varphi\colon\mathbb{G}_m\to\check K$ 为
\begin{equation}
    \check\varphi(\alpha)=\left(\alpha^{h\check\varphi_1},\ldots,\alpha^{h\check\varphi_n}\right) \tag{10.32}
\end{equation}
% [Tutor Note: 原文 (10.32) 的指数作 h（无 caron），而 (10.30) 定义的是 ǩ、且 (10.31) 有 ǎ_i = ǩǔ_i；此处照原文保留，请审校确认是否应为 ǩ。]
则它是单射。通过正合列 (10.18) 的转置
\begin{equation}
    1\to\check G_{\max}\to\check K\to\mathbb{G}_m\to 1 \tag{10.33}
\end{equation}
把 $\check G_{\max}$ 与 $\check K$ 的子群等同，则 $\operatorname{Im}\check\varphi\cap\check G_{\max}$ 由
\begin{equation}
    \check J:=\left(\exp\left(2\pi\sqrt{-1}\check\varphi_1\right),\ldots,\exp\left(2\pi\sqrt{-1}\check\varphi_n\right)\right) \tag{10.34}
\end{equation}
生成。记
\begin{equation}
    \check G:=\check G_{\max}\cap\mathrm{SL}_n(\C) \tag{10.35}
\end{equation}
则
\begin{equation}
    \overline{\check G}:=\check G/\langle\check J\rangle \tag{10.36}
\end{equation}
自然作用于 $\check Y$，故记对应的轨形为 $[\check Y/\overline{\check G}]$。像这样作出 Calabi–Yau 轨形的组 $(\mathbb{Y},[\check Y/\overline{\check G}])$ 的手续，称为转置镜像构造（transposition mirror construction）；它是 Greene–Plesser [GP90] 的轨形镜像构造（orbifold mirror construction）（的特殊情形）的推广，由 Berglund–Hübsch [BH93] 引入。

\begin{theorem}[{[CR11, Theorem 2]}]
对满足 Calabi–Yau 条件 (10.9) 的任意可逆多项式 $f$，$\mathbb{Y}$ 与 $[\check Y/\overline{\check G}]$ 是拓扑镜像对。
\end{theorem}

这里 $\mathbb{Y}$ 与 $[\check Y/\overline{\check G}]$ 是轨形，故拓扑镜像对的定义中使用的不是通常的 Hodge 数，而是后述的轨形 Hodge 数（orbifold Hodge number）。

[CR11, Theorem 2] 更强：它证明了把转置镜像构造按 [BH95, Kra] 推广到含群作用的形式后，仍给出拓扑镜像对。

\begin{theorem}[{[CR11, Theorem 2]}]
对满足 Calabi–Yau 条件 (10.9) 的任意可逆多项式 $f$，以及极大对角对称群 $G_{\max}$ 的满足 $\langle J\rangle\subset G\subset\mathrm{SL}_n(\C)$ 的子群 $G$，把 $G$ 的转置定义为
\begin{equation}
    \check G=\left\{\prod_{i=1}^{n}\check\rho_i^{r_i}\ \middle|\ (r_1\ \cdots\ r_n)A^{-1}
    \begin{pmatrix}
        a_1 \\ \vdots \\ a_n
    \end{pmatrix}
    \in\Z\ \text{ for all }\ \prod_{i=1}^{n}\rho_i^{a_i}\in G\right\} \tag{10.37}
\end{equation}
则轨形 $[\mathbb{Y}/\overline{G}]$ 与 $[\check Y/\overline{\check G}]$ 是拓扑镜像对。
\end{theorem}

这里 $\overline{G}$ 与 $\overline{\check G}$ 由下式定义。
\begin{equation}
    \overline{G}:=G/\langle J\rangle,\qquad \overline{\check G}:=\check G/\langle\check J\rangle \tag{10.38}
\end{equation}

% CHUNK c06 DONE pages p022–p026
